%% ============================================================
%%  Radar Target Classification: Classical ML vs Quantum ML
%%  A Study on Terrestrial Coastal Radar Data
%% ============================================================
\documentclass[12pt, a4paper]{article}

\usepackage[margin=2.5cm]{geometry}
\usepackage{amsmath, amssymb, amsfonts}
\usepackage{graphicx}
\usepackage{booktabs}
\usepackage{array}
\usepackage{multirow}
\usepackage{xcolor}
\usepackage{hyperref}
\usepackage{caption}
\usepackage{subcaption}
\usepackage{algorithm}
\usepackage{algpseudocode}
\usepackage{listings}
\usepackage{tikz}
\usepackage{pgfplots}
\pgfplotsset{compat=1.18}
\usepackage{float}
\usepackage{enumitem}
\usepackage[numbers]{natbib}
\usepackage{setspace}
\onehalfspacing

\hypersetup{
    colorlinks=true, linkcolor=blue, citecolor=blue, urlcolor=blue
}

\lstset{
    basicstyle=\ttfamily\small,
    breaklines=true,
    frame=single,
    backgroundcolor=\color{gray!10}
}

% ── Custom commands ──────────────────────────────────────────
\newcommand{\vect}[1]{\boldsymbol{#1}}
\newcommand{\mat}[1]{\mathbf{#1}}
\newcommand{\ket}[1]{|#1\rangle}
\newcommand{\bra}[1]{\langle #1|}
\newcommand{\braket}[2]{\langle #1 | #2 \rangle}
\newcommand{\expect}[1]{\langle #1 \rangle}
\DeclareMathOperator{\argmax}{arg\,max}

%% ============================================================
\begin{document}

\title{
    \textbf{Multi-Class Vessel Classification from Terrestrial Coastal Radar:\\
    A Comparative Study of Classical Machine Learning\\
    and Quantum Machine Learning Approaches}\\[0.5em]
    \large Working Paper — Feature Analysis, Methodology and Results
}
\author{
    [Author Name]\\
    \textit{[Institution / Affiliation]}\\
    \texttt{[email]}
}
\date{\today}
\maketitle

\begin{abstract}
This paper presents a comprehensive study of multi-class vessel type classification
using features extracted from a terrestrial X-band coastal surveillance radar.
Five vessel types — Fishing (Type 30), Dredger/Underwater Ops (Type 33),
Tug (Type 52), Cargo (Type 70), and Tanker (Type 80) — are distinguished
using radar-derived features including amplitude, spatial extent, and azimuth.
We systematically examine each feature's physical interpretation and statistical
separability, then compare the classification performance of classical machine
learning algorithms (Random Forest, Support Vector Machine, Multi-Layer
Perceptron, Gradient Boosting Trees) against three Quantum Machine Learning
(QML) approaches: Variational Quantum Classifier (VQC), Hybrid Quantum--Classical
Neural Network (HQNN), and Quantum Kernel Estimation (QKE).
A key contribution is the derivation of radar-native physical size descriptors
— down-range extent and cross-range physical extent — as discriminating features
for the traditionally difficult cargo vs.\ tanker boundary.
Results are evaluated by accuracy, macro F1-score, area under the ROC curve (AUC),
and model confidence score distributions.
\end{abstract}

\tableofcontents
\newpage

%% ============================================================
\section{Introduction}
\label{sec:introduction}

Terrestrial coastal radar systems continuously monitor maritime traffic by
detecting, tracking, and reporting the position and motion of vessels in their
surveillance zone. While radar provides reliable detection of physical objects,
identifying the \emph{type} of vessel from radar returns alone is a non-trivial
classification problem. Traditionally, vessel identity is obtained from the
Automatic Identification System (AIS), a transponder-based reporting protocol.
However, AIS can be switched off, spoofed, or simply absent on non-SOLAS vessels.
A radar system capable of inferring vessel type purely from its own measurements
would therefore provide an independent, passive classification capability of
significant maritime domain awareness value.

This paper addresses the following research questions:
\begin{enumerate}[label=\textbf{RQ\arabic*.}]
    \item Which radar-extracted features carry the most discriminating information
          for vessel type classification?
    \item Do classical ML algorithms achieve satisfactory classification accuracy
          on this naturally imbalanced, feature-overlapping dataset?
    \item Does Quantum Machine Learning offer measurable advantages over classical
          approaches in terms of accuracy, generalisation, or confidence calibration?
    \item Can the historically difficult Cargo (Type~70) vs.\ Tanker (Type~80)
          boundary be resolved by correctly derived radar-native size features?
\end{enumerate}

The paper is structured as follows.
Section~\ref{sec:radar_background} explains the physics of radar measurement and
how each feature is extracted.
Section~\ref{sec:dataset} describes the dataset, class distribution, and
statistical separability analysis.
Section~\ref{sec:preprocessing} covers data cleaning, normalisation, and class
balancing.
Section~\ref{sec:classical_ml} presents the classical ML models.
Section~\ref{sec:qml} introduces the QML approaches.
Section~\ref{sec:experiments} describes the experimental setup.
Section~\ref{sec:results} presents and compares the results.
Section~\ref{sec:discussion} discusses findings, confidence calibration, and
limitations.
Section~\ref{sec:conclusion} concludes the paper.

%% ============================================================
\section{Radar Background and Feature Extraction}
\label{sec:radar_background}

\subsection{Terrestrial Radar Operation}

A terrestrial coastal surveillance radar emits short pulses of electromagnetic
energy and listens for the reflections (echoes) returned by objects within its
coverage area. The fundamental measured quantities are:
\begin{itemize}
    \item \textbf{Range} $r$: the radial distance from the radar antenna to the
          target, inferred from the round-trip travel time of the pulse,
          $r = c \cdot \Delta t / 2$, where $c$ is the speed of light.
    \item \textbf{Azimuth} $\phi$: the angular bearing of the target measured
          clockwise from North, determined by the orientation of the rotating
          antenna at the moment of detection.
    \item \textbf{Amplitude}: the strength of the received echo, related to the
          target's Radar Cross Section (RCS).
\end{itemize}

A vessel occupies multiple range and azimuth resolution cells simultaneously.
The radar processing pipeline groups contiguous cells that exceed a detection
threshold into a single \emph{plot}, and a tracking algorithm associates
successive plots into a \emph{track} (the vessel's trajectory over time).

\subsection{Feature Dictionary}
\label{sec:features}

Table~\ref{tab:features} summarises all extracted features.
The subsections below explain each feature in detail.

\begin{table}[H]
\centering
\caption{Complete feature dictionary. All features are radar-native unless noted.}
\label{tab:features}
\begin{tabular}{llp{7cm}}
\toprule
\textbf{Feature} & \textbf{Unit} & \textbf{Description} \\
\midrule
\multicolumn{3}{l}{\textit{Spatial position}} \\
\texttt{range}              & m           & Radial distance from antenna to target \\
\texttt{azimuth}            & rad         & Angular bearing (0 = North, clockwise) \\
\texttt{RLatitude}          & deg         & Estimated target latitude \\
\texttt{RLongitude}         & deg         & Estimated target longitude \\
\midrule
\multicolumn{3}{l}{\textit{Amplitude / Radar Cross Section}} \\
\texttt{PeakAmplitude}      & linear      & Maximum received power across the plot window \\
\texttt{TotalAmplitude}     & linear      & Sum of received power across all cells in the plot window \\
\texttt{PeakAmplitude\_corr}  & linear    & Range-corrected peak amplitude (see Eq.~\ref{eq:amp_corr}) \\
\texttt{TotalAmplitude\_corr} & linear    & Range-corrected total amplitude \\
\midrule
\multicolumn{3}{l}{\textit{Spatial extent — raw gate widths}} \\
\texttt{winrgw}             & m           & Window range gate width (full detection window in range) \\
\texttt{winazw}             & rad         & Window azimuth gate width (full detection window in azimuth) \\
\texttt{rgw} / \texttt{down\_range\_extent}  & m  & Range gate width of the plot (vessel extent along range axis) \\
\texttt{azw} / \texttt{cross\_range\_extent} & rad & Azimuth gate width of the plot (vessel extent along azimuth axis) \\
\midrule
\multicolumn{3}{l}{\textit{Spatial extent — derived physical dimensions}} \\
\texttt{az\_extent\_m}      & m           & Physical azimuth extent: $r \times \texttt{azw}$ (Eq.~\ref{eq:az_extent}) \\
\texttt{euclid\_size}       & m           & Euclidean size: $\sqrt{\texttt{rgw}^2 + \texttt{azw}^2} \approx \texttt{rgw}$ \\
\texttt{ellipse\_area}      & m$^2$       & $\pi \cdot (\texttt{rgw}/2) \cdot (\texttt{azw}/2)$ \\
\texttt{size\_beam\_component}     & m   & Size projected onto vessel beam axis (aspect-dependent) \\
\texttt{size\_bow\_stern\_component} & m & Size projected onto vessel bow–stern axis (aspect-dependent) \\
\texttt{cr\_dr\_ratio}      & —           & $\texttt{azw} / \texttt{rgw}$: shape elongation ratio \\
\midrule
\multicolumn{3}{l}{\textit{Motion}} \\
\texttt{RSog}               & kn          & Radar-measured speed over ground \\
\texttt{RCog}               & rad         & Radar-measured course over ground \\
\texttt{AvgRSog}            & kn          & Moving average of radar speed \\
\texttt{aspect}             & rad         & Aspect angle: angle between vessel heading and radar line-of-sight \\
\texttt{aspect\_fold}       & rad         & Folded aspect (symmetrised, range $[0, \pi/2]$) \\
\texttt{aspect\_sin}        & —           & $\sin(\texttt{aspect\_fold})$ \\
\texttt{aspect\_cos}        & —           & $\cos(\texttt{aspect\_fold})$ \\
\midrule
\multicolumn{3}{l}{\textit{Temporal rolling statistics (window $w \in \{900, 1800, 3600, 10800\}$ s)}} \\
\texttt{measured\_sog\_avg\_$w$}    & kn  & Mean speed over window $w$ \\
\texttt{measured\_sog\_std\_$w$}    & kn  & Speed variability over window $w$ \\
\texttt{measured\_cog\_avg\_$w$}    & rad & Mean course over window $w$ \\
\texttt{measured\_cog\_std\_$w$}    & rad & Course variability (heading steadiness) \\
\texttt{measured\_size\_avg\_$w$}   & m   & Mean track size over window $w$ \\
\texttt{measured\_rangeStd\_$w$}    & m   & Std of range positions (manoeuvring indicator) \\
\texttt{measured\_azimuthStd\_$w$}  & rad & Std of azimuth positions \\
\midrule
\multicolumn{3}{l}{\textit{AIS-derived labels (not used as model features)}} \\
\texttt{AIS\_LENGTH}        & m           & Vessel physical length from AIS transponder \\
\texttt{AIS\_WIDTH}         & m           & Vessel physical beam from AIS transponder \\
\texttt{Type}               & —           & AIS vessel type code (ground-truth label) \\
\texttt{size\_class}        & —           & Categorical: small / medium / large \\
\bottomrule
\end{tabular}
\end{table}

\subsubsection{Azimuth ($\phi$)}

The azimuth is the horizontal bearing angle from the radar antenna to the target,
measured in radians (or degrees) clockwise from North:
\begin{equation}
    \phi = \text{atan2}(\Delta E,\, \Delta N) \mod 2\pi
\end{equation}
where $\Delta E$ and $\Delta N$ are the East and North displacements.
For a rotating antenna, azimuth is directly read from the antenna encoder at
the moment of detection. The azimuth value for a given vessel varies as the
vessel moves around the radar coverage area, making it a \emph{positional}
rather than a vessel-intrinsic feature. Nevertheless, different vessel types
concentrate in different operational zones (e.g., tankers in deep-water channels,
fishing vessels in shallower sectors), giving azimuth moderate discriminating
power (Cohen's $d \approx 0.51$ for Cargo vs.\ Tanker in our dataset).

\subsubsection{Amplitude Features}

\paragraph{Raw amplitudes.}
The radar receiver records the signal power in each range-azimuth resolution cell.
For a detection, two aggregate amplitude values are computed over the plot window:

\begin{align}
    \texttt{PeakAmplitude} &= \max_{i \in \text{plot}} A_i \label{eq:peak}\\
    \texttt{TotalAmplitude} &= \sum_{i \in \text{plot}} A_i \label{eq:total}
\end{align}

where $A_i$ is the received power amplitude in cell $i$.
\texttt{PeakAmplitude} captures the brightest point on the vessel (often the
bridge or superstructure). \texttt{TotalAmplitude} captures the integrated RCS
over the entire vessel footprint and is proportional to the vessel's total
reflecting area.

\paragraph{Range-corrected amplitudes.}
Raw radar amplitude decays with range according to the radar range equation.
For monostatic radar under free-space propagation:
\begin{equation}
    P_r = \frac{P_t G^2 \lambda^2 \sigma}{(4\pi)^3 r^4}
    \label{eq:radar_eq}
\end{equation}
where $P_r$ is received power, $P_t$ is transmitted power, $G$ is antenna gain,
$\lambda$ is wavelength, $\sigma$ is the target RCS, and $r$ is range.
To make amplitude features range-invariant (i.e., reflecting the vessel's
\emph{intrinsic} RCS rather than its distance), range-corrected values are
computed:
\begin{equation}
    \texttt{PeakAmplitude\_corr} = \texttt{PeakAmplitude} \cdot r^4 / C
    \label{eq:amp_corr}
\end{equation}
where $C$ is a system constant. Range-corrected amplitudes are more physically
meaningful for comparing vessels at different ranges.

\paragraph{Physical interpretation for vessel classification.}
Larger vessels (Tanker, Cargo) generally exhibit higher \texttt{TotalAmplitude}
than smaller vessels (Tug, Fishing), but within large-vessel classes the
amplitude overlap is significant because RCS depends not only on size but also
on aspect angle, hull material, and sea state.

\subsubsection{Spatial Extent Features}
\label{sec:size_features}

A vessel occupies a finite number of range and azimuth resolution cells.
The \emph{extent} of the detection in each axis provides a measure of the
vessel's physical dimensions as seen by the radar.

\paragraph{Range gate width (\texttt{rgw} / \texttt{down\_range\_extent}).}
This is the number of range resolution cells occupied by the detection, converted
to metres by multiplying by the range resolution $\delta r$:
\begin{equation}
    \texttt{rgw} = N_r \cdot \delta r \quad [\text{metres}]
\end{equation}
It approximates the vessel's \emph{length projected along the radar line-of-sight}
(the down-range dimension). For a vessel at aspect angle $\alpha$ with true
length $L$:
\begin{equation}
    \texttt{rgw} \approx L \cos\alpha + W \sin\alpha
\end{equation}
where $W$ is the vessel beam width.

\paragraph{Azimuth gate width (\texttt{azw} / \texttt{cross\_range\_extent}).}
This is the angular span of the detection in the azimuth dimension, in radians:
\begin{equation}
    \texttt{azw} = N_\phi \cdot \delta\phi \quad [\text{radians}]
\end{equation}
where $N_\phi$ is the number of azimuth cells and $\delta\phi$ is the azimuth
resolution. Note that \texttt{azw} is in \emph{angular} units; it must be
converted to physical length before size comparison.

\paragraph{Physical azimuth extent (\texttt{az\_extent\_m}).}
The arc length subtended by the detection at range $r$ is:
\begin{equation}
    \texttt{az\_extent\_m} = r \cdot \texttt{azw} \quad [\text{metres}]
    \label{eq:az_extent}
\end{equation}
This is the vessel's \emph{cross-range physical extent} — the dimension
perpendicular to the radar beam direction. For a vessel at aspect $\alpha$:
\begin{equation}
    \texttt{az\_extent\_m} \approx L \sin\alpha + W \cos\alpha
\end{equation}
Together, \texttt{rgw} and \texttt{az\_extent\_m} represent the two orthogonal
projections of the vessel's hull geometry onto the radar measurement axes.
A tanker, being wide relative to its length (low length-to-beam ratio), tends
to show a larger \texttt{az\_extent\_m} than a cargo vessel of similar length.
In our dataset, mean \texttt{az\_extent\_m} for Tanker (Type~80) is 898~m
versus 634~m for Cargo (Type~70), giving Cohen's $d = 0.502$.

\paragraph{Why \texttt{euclid\_size} $\approx$ \texttt{rgw} in this dataset.}
A common derived size metric is:
\begin{equation}
    \texttt{euclid\_size} = \sqrt{\texttt{rgw}^2 + \texttt{azw}^2}
\end{equation}
However, because \texttt{azw} is stored in radians ($\sim$0.02), while
\texttt{rgw} is in metres ($\sim$200), the azimuth term is negligible:
\begin{equation}
    \sqrt{200^2 + 0.02^2} \approx 200 \approx \texttt{rgw}
\end{equation}
Therefore \texttt{euclid\_size} is \emph{numerically identical} to \texttt{rgw}
in this dataset and carries no additional information. The same applies to
\texttt{ellipse\_area}:
\begin{equation}
    \texttt{ellipse\_area} = \pi \cdot \frac{\texttt{rgw}}{2} \cdot \frac{\texttt{azw}}{2}
    \approx \pi \cdot 100 \cdot 0.01 \approx 3.1
\end{equation}
This near-zero value is meaningless as a discriminating feature.
The physically correct formulation requires converting azimuth to metres first
(i.e., using \texttt{az\_extent\_m}), then computing:
\begin{align}
    \texttt{euclid\_size\_phys} &= \sqrt{\texttt{rgw}^2 + \texttt{az\_extent\_m}^2} \\
    \texttt{ellipse\_area\_phys} &= \pi \cdot \frac{\texttt{rgw}}{2} \cdot \frac{\texttt{az\_extent\_m}}{2}
\end{align}

\paragraph{Aspect-decomposed size components.}
Given the estimated aspect angle $\alpha$ (angle between vessel heading and
radar line-of-sight), the size can be decomposed into:
\begin{align}
    \texttt{size\_bow\_stern\_component} &= \texttt{euclid\_size} \cdot \cos\alpha \\
    \texttt{size\_beam\_component}       &= \texttt{euclid\_size} \cdot \sin\alpha
\end{align}
When the vessel is bow-on ($\alpha \approx 0$), the bow--stern component
dominates; broadside ($\alpha \approx \pi/2$) gives a large beam component.

\subsubsection{Motion and Aspect Features}

\paragraph{Radar speed and course (\texttt{RSog}, \texttt{RCog}).}
The radar tracker estimates vessel velocity from successive plot positions:
\begin{equation}
    \texttt{RSog} = \frac{\|\vect{p}_{t} - \vect{p}_{t-1}\|}{\Delta t}, \qquad
    \texttt{RCog} = \text{atan2}(p_{E,t}-p_{E,t-1},\, p_{N,t}-p_{N,t-1})
\end{equation}
Speed and course provide behavioural signatures: tankers typically cruise at
steady 10--14~knots on fixed routes; fishing vessels show low, variable speeds
with frequent course changes.

\paragraph{Aspect angle ($\alpha$).}
The aspect angle is the angle between the vessel's heading $\psi$ and the
direction from the vessel to the radar $\phi_{\text{to-radar}}$:
\begin{equation}
    \alpha = |\psi - \phi_{\text{to-radar}}|
\end{equation}
Because the radar signature of a vessel is highly aspect-dependent (a broadside
vessel appears much larger than a bow-on vessel), aspect is important for
interpreting amplitude and size features. Three representations are used:
\begin{align}
    \texttt{aspect\_fold}  &= \min(\alpha,\, \pi - \alpha) \in [0,\, \pi/2] \\
    \texttt{aspect\_sin}   &= \sin(\texttt{aspect\_fold}) \\
    \texttt{aspect\_cos}   &= \cos(\texttt{aspect\_fold})
\end{align}
The folded representation removes the bow/stern ambiguity (a vessel looks the
same from the front and the back).

\subsubsection{Temporal Rolling Statistics}

For each vessel track, rolling statistics are computed over four time windows:
$w \in \{900, 1800, 3600, 10800\}$ seconds (15~min, 30~min, 1~h, 3~h).
For each window $w$ ending at time $t$, let $\mathcal{W} = \{s : t-w \le s \le t\}$:

\begin{align}
    \texttt{measured\_sog\_avg\_}w &= \frac{1}{|\mathcal{W}|}\sum_{s \in \mathcal{W}} \text{SOG}_s \\
    \texttt{measured\_sog\_std\_}w &= \text{std}_{s \in \mathcal{W}}(\text{SOG}_s) \\
    \texttt{measured\_cog\_std\_}w &= \text{std}_{s \in \mathcal{W}}(\text{COG}_s) \quad \text{(heading steadiness)}\\
    \texttt{measured\_rangeStd\_}w &= \text{std}_{s \in \mathcal{W}}(r_s) \quad \text{(spatial spread in range)}
\end{align}

These temporal features capture \emph{behavioural patterns}:
\begin{itemize}
    \item A low \texttt{measured\_cog\_std} indicates a vessel maintaining a
          steady course (characteristic of large transit vessels: Cargo, Tanker).
    \item High \texttt{measured\_sog\_std} indicates speed variability
          (characteristic of Fishing, Tug).
    \item Low \texttt{measured\_rangeStd} for a 3-hour window indicates
          stationary or slow-moving operation (Dredger, anchor).
\end{itemize}

%% ============================================================
\section{Dataset Description}
\label{sec:dataset}

\subsection{Data Collection}

The dataset was collected from a network of X-band terrestrial coastal radar
stations operating in a coastal surveillance environment. Raw radar tracks
(plots) are fused with AIS reports to provide ground-truth vessel type labels
using the ITU-R AIS vessel type codes. After AIS correlation, each radar
observation carries a confirmed vessel type label.

\subsection{Vessel Type Classes}

Five vessel types are selected for classification based on their operational
relevance and AIS type code prevalence:

\begin{table}[H]
\centering
\caption{Target classification classes.}
\label{tab:classes}
\begin{tabular}{cll}
\toprule
\textbf{Type Code} & \textbf{Vessel Category} & \textbf{Characteristics} \\
\midrule
30 & Fishing           & Small to medium, variable speed, non-linear tracks \\
33 & Dredging / Underwater Ops & Medium to large, very low speed or stationary \\
52 & Tug               & Small, powerful, assists other vessels \\
70 & Cargo             & Large, steady speed 10--14~kn, straight tracks \\
80 & Tanker            & Large, slightly wider beam than cargo, steady transit \\
\bottomrule
\end{tabular}
\end{table}

\subsection{Dataset Statistics}

\begin{table}[H]
\centering
\caption{Dataset statistics after cleaning (zero-amplitude rows removed).}
\label{tab:dataset_stats}
\begin{tabular}{lrrrrrr}
\toprule
\textbf{Feature} & \textbf{Type 30} & \textbf{Type 33} & \textbf{Type 52}
                 & \textbf{Type 70} & \textbf{Type 80} & \textbf{Overall} \\
\midrule
Count           & 3{,}729  & 6{,}022  & 3{,}540  & 59{,}041  & 28{,}098  & 100{,}430 \\
\midrule
azimuth (mean)          & --   & --  & --  & 2.3 & 3.0 & -- \\
PeakAmplitude (mean)    & --   & --  & --  & 182 & 174 & -- \\
TotalAmplitude (mean)   & --   & --  & -- & 85{,}824 & 103{,}890 & -- \\
down\_range\_ext (mean) & 169  & 353 & 188 & 293 & 270 & 240 \\
az\_extent\_m (mean)    & 262  & 590 & 499 & 634 & 898 & 585 \\
\bottomrule
\end{tabular}
\end{table}

\subsection{Class Imbalance}

The dataset exhibits significant class imbalance, with Cargo (Type~70) accounting
for 59\% of all observations. The class ratio between the largest (Type~70:
59,041) and smallest (Type~30: 3,729) class is approximately 16:1.
This imbalance must be addressed during training (see Section~\ref{sec:balancing}).

\subsection{Feature Separability Analysis}
\label{sec:separability}

Cohen's $d$ measures the standardised mean difference between two classes:
\begin{equation}
    d = \frac{\mu_A - \mu_B}{\sqrt{(\sigma_A^2 + \sigma_B^2)/2}}
\end{equation}
A value of $|d| < 0.5$ indicates poor separability, $0.5$--$0.8$ moderate,
and $>0.8$ strong. Table~\ref{tab:cohen} shows Cohen's $d$ for the critical
Cargo (70) vs.\ Tanker (80) boundary under different feature sets.

\begin{table}[H]
\centering
\caption{Cohen's $d$ for Cargo (Type~70) vs.\ Tanker (Type~80).}
\label{tab:cohen}
\begin{tabular}{lrll}
\toprule
\textbf{Feature} & \textbf{Cohen's $d$} & \textbf{Interpretation} & \textbf{Dataset} \\
\midrule
azimuth             & 0.508 & Moderate & radarfeatureL\_Study \\
az\_extent\_m       & 0.502 & Moderate & radarfeatureL\_Study \\
TotalAmplitude      & 0.231 & Poor     & radarfeatureL\_Study \\
down\_range\_extent & 0.240 & Poor     & radarfeatureL\_Study \\
PeakAmplitude       & 0.203 & Poor     & radarfeatureL\_Study \\
\midrule
azimuth (original 5-feat) & 0.454 & Poor & Radar data.csv \\
TotalAmplitude            & 0.312 & Poor & Radar data.csv \\
ellipse\_area             & 0.361 & Poor & Radar data.csv \\
euclid\_size              & 0.320 & Poor & Radar data.csv \\
PeakAmplitude             & 0.191 & Poor & Radar data.csv \\
\bottomrule
\end{tabular}
\end{table}

The introduction of \texttt{az\_extent\_m} as a physically correct cross-range
extent (in metres) brings the 70/80 Cohen's $d$ from below 0.5 to 0.502,
the first time any feature in this study achieves separability at or above
the moderate threshold for this class pair.

%% ============================================================
\section{Data Preprocessing}
\label{sec:preprocessing}

\subsection{Cleaning}
\label{sec:cleaning}

The following cleaning steps are applied:
\begin{enumerate}
    \item \textbf{Zero-amplitude removal}: Rows where \texttt{PeakAmplitude}~$= 0$
          are removed. These correspond to ghost tracks or failed plot associations
          where no actual radar return was linked to the track record.
          In the study dataset, 17,532 of 117,962 target-class rows (14.9\%) are
          removed by this filter.
    \item \textbf{Class filtering}: Only the five target classes
          (30, 33, 52, 70, 80) are retained.
    \item \textbf{Feature derivation}: \texttt{az\_extent\_m} is computed as
          $r \times \texttt{cross\_range\_extent}$.
    \item \textbf{Null removal}: Rows with null values in any selected feature
          column are dropped.
\end{enumerate}

\subsection{Feature Normalisation}
\label{sec:normalisation}

All features are standardised using the training-set mean $\mu$ and standard
deviation $\sigma$ (StandardScaler):
\begin{equation}
    x' = \frac{x - \mu}{\sigma}
\end{equation}
The scaler parameters are fitted on the training set only and applied
identically to the test set and any new inference data, preventing data leakage.

\subsection{Class Balancing with SMOTE}
\label{sec:balancing}

The Synthetic Minority Over-sampling Technique (SMOTE)~\cite{chawla2002smote}
is applied to the training set to address class imbalance. SMOTE generates
synthetic samples for minority classes by interpolating between existing
minority-class samples in feature space:
\begin{equation}
    \vect{x}_{\text{new}} = \vect{x}_i + \lambda \cdot (\vect{x}_{k} - \vect{x}_i),
    \quad \lambda \sim \mathcal{U}(0,1)
\end{equation}
where $\vect{x}_i$ is an existing minority sample and $\vect{x}_k$ is one of
its $K$ nearest neighbours ($K=5$ by default). SMOTE is applied exclusively
to the training set; the test set retains the original class distribution.

Additionally, \textbf{class-weighted loss} is used during neural network and
quantum model training:
\begin{equation}
    w_c = \frac{N}{n_c \cdot C}
\end{equation}
where $N$ is total training samples, $n_c$ is the count of class $c$, and
$C$ is the number of classes.

\subsection{Train / Test Split}

An 80/20 stratified split is used, preserving the original class proportions
in both subsets. The split is performed \emph{before} SMOTE to ensure no
synthetic samples appear in the test set.

%% ============================================================
\section{Classical Machine Learning Models}
\label{sec:classical_ml}

\subsection{Random Forest (RF)}

Random Forest~\cite{breiman2001random} builds an ensemble of $T$ decision trees,
each trained on a bootstrapped sample and a random subset of features. The
predicted class is the majority vote across all trees:
\begin{equation}
    \hat{y} = \text{mode}\{h_t(\vect{x})\}_{t=1}^T
\end{equation}
RF is robust to overfitting, handles mixed-scale features without normalisation,
and provides feature importance estimates. It serves as the baseline classical
model in this study.

\subsection{Support Vector Machine (SVM)}

SVM~\cite{cortes1995support} finds the maximum-margin hyperplane separating
classes in a (possibly kernelised) feature space. For multi-class problems,
a one-vs-one scheme is used. With the Radial Basis Function (RBF) kernel:
\begin{equation}
    K(\vect{x}_i, \vect{x}_j) = \exp\!\left(-\gamma \|\vect{x}_i - \vect{x}_j\|^2\right)
\end{equation}
SVM is effective in high-dimensional spaces but computationally expensive
at scale ($\mathcal{O}(N^2)$--$\mathcal{O}(N^3)$).

\subsection{Multi-Layer Perceptron (MLP)}

A feedforward neural network with $L$ hidden layers:
\begin{equation}
    \vect{h}^{(l)} = \sigma\!\left(\mat{W}^{(l)} \vect{h}^{(l-1)} + \vect{b}^{(l)}\right),
    \quad \vect{h}^{(0)} = \vect{x}
\end{equation}
where $\sigma$ is the ReLU activation. The output layer applies softmax to
produce class probabilities. Trained with Adam optimiser and cross-entropy loss.

\subsection{Gradient Boosting Trees (GBT)}

GBT~\cite{friedman2001greedy} builds an additive ensemble of weak learners
(shallow decision trees) in a stage-wise fashion, where each tree corrects
the residual errors of its predecessors:
\begin{equation}
    F_m(\vect{x}) = F_{m-1}(\vect{x}) + \eta \cdot h_m(\vect{x})
\end{equation}
where $\eta$ is the learning rate and $h_m$ is fitted to the negative gradient
of the loss. GBT excels at capturing non-linear feature interactions
(e.g., large \texttt{az\_extent\_m} \emph{and} low speed $\Rightarrow$ Tanker)
that are difficult for linear or single-feature-based models.

%% ============================================================
\section{Quantum Machine Learning Models}
\label{sec:qml}

\subsection{Quantum Computing Primer}

A quantum computer operates on \emph{qubits} — quantum analogues of classical
bits. A single qubit state is a superposition:
\begin{equation}
    \ket{\psi} = \alpha\ket{0} + \beta\ket{1}, \quad |\alpha|^2 + |\beta|^2 = 1
\end{equation}
An $n$-qubit system lives in a $2^n$-dimensional Hilbert space, enabling
exponentially large state spaces. Quantum gates are unitary operations $U$
acting on one or more qubits. A quantum circuit is a sequence of gates applied
to an initial state $\ket{0}^{\otimes n}$.

\subsection{Data Encoding}

Classical feature vectors must be encoded into quantum states before processing.
We use \textbf{angle encoding} (also called \emph{rotational encoding}):
\begin{equation}
    \text{Encode}(\vect{x}) = \bigotimes_{i=1}^{n} R_y(x_i) \ket{0}
\end{equation}
where $R_y(\theta) = \exp(-i\theta Y/2)$ is a rotation about the $Y$-axis
by angle $\theta$. One qubit encodes one feature; for $d$ features, $d$ qubits
are required. Features are normalised to $[0, \pi]$ before encoding.
After encoding, qubit $i$ holds:
\begin{equation}
    R_y(x_i)\ket{0} = \cos\!\tfrac{x_i}{2}\ket{0} + \sin\!\tfrac{x_i}{2}\ket{1}
\end{equation}

\subsection{Variational Quantum Classifier (VQC)}
\label{sec:vqc}

The VQC~\cite{schuld2020circuit} is a parameterised quantum circuit used
as a classifier. It consists of three components:

\begin{enumerate}
    \item \textbf{Feature embedding layer} $U_\phi(\vect{x})$: encodes input
          features into the quantum state using angle encoding.
    \item \textbf{Variational (ansatz) layer} $U_\theta(\vect{\theta})$: a
          sequence of parameterised single-qubit rotations and two-qubit
          entangling gates, repeated for $L$ layers.
    \item \textbf{Measurement layer}: Pauli-$Z$ expectation values
          $\langle Z_i \rangle = \bra{\psi}Z_i\ket{\psi} \in [-1, 1]$
          are measured on each qubit and passed through a classical linear
          layer to produce class logits.
\end{enumerate}

We use the \textbf{Strongly Entangling Layers} ansatz~\cite{pennylane2019},
which applies $R_z$, $R_y$, $R_z$ rotations on each qubit followed by a
ring of CNOT gates, creating entanglement between all qubits:
\begin{equation}
    U_\text{SEL}(\vect{\theta}) = \prod_{l=1}^{L} \left[
        \prod_{i=1}^{n} R_z(\theta_{l,i,1}) R_y(\theta_{l,i,2}) R_z(\theta_{l,i,3})
        \cdot \text{CNOT}_{i \to i+1}
    \right]
\end{equation}

The full VQC circuit is:
\begin{equation}
    \ket{\psi(\vect{x}, \vect{\theta})} = U_\theta(\vect{\theta})\, U_\phi(\vect{x}) \ket{0}^{\otimes n}
\end{equation}

The output class probabilities are:
\begin{equation}
    p_c = \text{softmax}\!\left(\mat{W}_{\text{out}} \cdot [\langle Z_1 \rangle, \dots, \langle Z_n \rangle]^\top + \vect{b}\right)_c
\end{equation}

\textbf{Configuration used:}
$n = 5$ qubits (one per feature), $L = 3$ ansatz layers,
Adam optimiser, learning rate $\eta = 0.005$, early stopping (patience = 20).

\subsection{Hybrid Quantum--Classical Neural Network (HQNN)}
\label{sec:hqnn}

The HQNN~\cite{mari2020transfer} combines classical and quantum layers:
\begin{enumerate}
    \item \textbf{Classical pre-processing layers}: Fully connected layers
          with ReLU activation transform the input features before quantum encoding.
          This allows the network to learn a better feature representation
          for the quantum circuit.
    \item \textbf{Quantum layer} (PennyLane TorchLayer): the VQC sub-circuit
          with angle encoding and strongly entangling ansatz.
    \item \textbf{Classical post-processing layers}: Further linear layers map
          quantum measurements to final class probabilities.
\end{enumerate}

\begin{equation}
    \hat{\vect{y}} = \text{softmax}\!\left(
        f_{\text{post}}\!\left(
            Q_\theta\!\left(f_{\text{pre}}(\vect{x})\right)
        \right)
    \right)
\end{equation}

\textbf{Configuration used:}
$n = 5$ qubits, $L = 2$ quantum layers, 2 classical hidden layers of width 32,
Adam, $\eta = 0.005$, early stopping (patience = 20).

\subsection{Quantum Kernel Estimation (QKE)}
\label{sec:qke}

QKE~\cite{havlicek2019supervised} computes a kernel matrix using the quantum
feature map $\phi : \vect{x} \mapsto \ket{\phi(\vect{x})}$, then trains a
classical Support Vector Machine on this kernel:
\begin{equation}
    k(\vect{x}_i, \vect{x}_j) = |\braket{\phi(\vect{x}_i)}{\phi(\vect{x}_j)}|^2
\end{equation}
The quantum advantage (if any) arises when the quantum feature map induces a
kernel that is exponentially hard to compute classically. QKE has quadratic
complexity $\mathcal{O}(N^2)$ in the number of training samples, making it
expensive for large datasets.

\subsection{Training with PennyLane and PyTorch}

All quantum models are implemented using PennyLane~\cite{pennylane2019} with
the PyTorch interface. The quantum circuit is embedded as a \texttt{TorchLayer},
making the full model differentiable and trainable with standard PyTorch
optimisers via the \emph{parameter-shift rule}:
\begin{equation}
    \frac{\partial}{\partial \theta_i} \langle O \rangle = \frac{1}{2}\left[
        \langle O \rangle_{\theta_i + \pi/2} - \langle O \rangle_{\theta_i - \pi/2}
    \right]
\end{equation}
This allows exact gradient computation without finite differences.

%% ============================================================
\section{Ensemble Methods}
\label{sec:ensemble}

\subsection{Soft-Vote Ensemble}

A soft-vote ensemble averages the class probability vectors from multiple models:
\begin{equation}
    p_c^{\text{ens}} = \frac{1}{M} \sum_{m=1}^{M} p_c^{(m)}(\vect{x})
\end{equation}
The predicted class is $\hat{y} = \argmax_c p_c^{\text{ens}}$.
By combining models with different inductive biases (e.g., quantum and classical),
the ensemble can reduce individual model errors.

\subsection{Confidence Score}

The confidence of a prediction is defined as the maximum class probability:
\begin{equation}
    \text{conf}(\vect{x}) = \max_c p_c^{\text{ens}}(\vect{x})
\end{equation}
A well-calibrated model should show higher confidence on correctly classified
samples than on misclassified ones. Comparing confidence distributions across
classical and quantum models provides insight into model calibration.

%% ============================================================
\section{Experimental Setup}
\label{sec:experiments}

\subsection{Hardware and Software}

\begin{itemize}
    \item \textbf{Hardware}: NVIDIA DGX workstation, multi-core CPU (quantum
          simulation is CPU-bound; no quantum hardware used).
    \item \textbf{Quantum simulator}: PennyLane \texttt{default.qubit}
          (statevector simulation).
    \item \textbf{Framework}: PyTorch 2.x + PennyLane 0.3x + scikit-learn.
    \item \textbf{Optimiser}: Adam with $\beta_1 = 0.9$, $\beta_2 = 0.999$.
\end{itemize}

\subsection{Hyperparameter Summary}

\begin{table}[H]
\centering
\caption{Hyperparameters for all models.}
\label{tab:hyperparams}
\begin{tabular}{llr}
\toprule
\textbf{Model} & \textbf{Parameter} & \textbf{Value} \\
\midrule
\multirow{4}{*}{VQC}  & Qubits   & 5 \\
                      & Layers   & 3 \\
                      & LR       & 0.005 \\
                      & Patience & 20 \\
\midrule
\multirow{5}{*}{HQNN} & Qubits         & 5 \\
                      & Quantum layers & 2 \\
                      & Classical hidden & 32 \\
                      & LR             & 0.005 \\
                      & Patience       & 20 \\
\midrule
\multirow{4}{*}{GBT}  & Estimators  & 300 \\
                      & Max depth   & 5 \\
                      & LR          & 0.05 \\
                      & Subsample   & 0.8 \\
\midrule
\multirow{2}{*}{RF}   & Estimators  & 300 \\
                      & Max features & \texttt{sqrt} \\
\midrule
SVM  & Kernel & RBF ($\gamma$ = scale) \\
\bottomrule
\end{tabular}
\end{table}

\subsection{Evaluation Metrics}

\begin{itemize}
    \item \textbf{Accuracy}: $\text{Acc} = \frac{\text{TP}_{\text{all}}}{N}$
    \item \textbf{Macro F1}: unweighted mean of per-class F1 scores,
          treating all classes equally regardless of support.
    \item \textbf{Macro AUC}: area under the one-vs-rest ROC curve, averaged
          across all classes.
    \item \textbf{Per-class Recall}: $R_c = \text{TP}_c / (\text{TP}_c + \text{FN}_c)$
    \item \textbf{Confidence distribution}: histogram of $\max_c p_c$ for correct
          and incorrect predictions.
\end{itemize}

%% ============================================================
\section{Results}
\label{sec:results}

\subsection{Model Comparison — 5-Feature Baseline (Original Dataset)}

\begin{table}[H]
\centering
\caption{Results on original 5 features (azimuth, TotalAmplitude, ellipse\_area,
euclid\_size, PeakAmplitude). Test set: 124 samples.}
\label{tab:results_5feat}
\begin{tabular}{lrrrr}
\toprule
\textbf{Model} & \textbf{Accuracy} & \textbf{F1 Macro} & \textbf{AUC} & \textbf{Best Epoch} \\
\midrule
VQC (lr=0.005, ES ep8)           & 55.6\% & 56.8\% & 0.780 & 8 \\
HQNN (lr=0.005, ES ep30)         & 66.9\% & 68.0\% & 0.878 & 30 \\
Ensemble (VQC + HQNN)            & 66.1\% & 68.1\% & 0.866 & -- \\
\midrule
\textit{(Classical baselines — to be completed)} & & & & \\
\bottomrule
\end{tabular}
\end{table}

\subsection{Model Comparison — Corrected 5-Feature Set (Study Dataset)}

\begin{table}[H]
\centering
\caption{Results on corrected 5 features (azimuth, PeakAmplitude, TotalAmplitude,
down\_range\_extent, az\_extent\_m). 10,000 training samples (2,000/class).}
\label{tab:results_5feat_new}
\begin{tabular}{lrrrr}
\toprule
\textbf{Model} & \textbf{Accuracy} & \textbf{F1 Macro} & \textbf{AUC} & \textbf{Notes} \\
\midrule
VQC                     & -- & -- & -- & \textit{(pending)} \\
HQNN                    & -- & -- & -- & \textit{(pending)} \\
GBT                     & -- & -- & -- & \textit{(pending)} \\
Ensemble (VQC+HQNN+GBT) & -- & -- & -- & \textit{(pending)} \\
\bottomrule
\end{tabular}
\end{table}

\subsection{Per-Class Recall}

\begin{table}[H]
\centering
\caption{Per-class recall (\%) — best model from each approach.}
\label{tab:per_class}
\begin{tabular}{lrrrrr}
\toprule
\textbf{Model} & \textbf{Type 30} & \textbf{Type 33} & \textbf{Type 52}
               & \textbf{Type 70} & \textbf{Type 80} \\
\midrule
HQNN (5-feat orig) & 100\% & 82\% & 71\% & 67\% & 47\% \\
Ensemble (merged)  & 98\%  & 89\% & 80\% & 56\% & 50\% \\
Ensemble (study)   & --    & --   & --   & --   & -- \\
\bottomrule
\end{tabular}
\end{table}

%% ============================================================
\section{Discussion}
\label{sec:discussion}

\subsection{The Cargo vs.\ Tanker Boundary}

The Cargo (Type~70) vs.\ Tanker (Type~80) confusion is the dominant classification
challenge in this study. All features from the original dataset had Cohen's
$d < 0.5$, placing this boundary below the threshold of statistical separability
by any individual feature. The introduction of the physically correct
cross-range extent \texttt{az\_extent\_m} = $r \times \texttt{azw}$ brings
the separability marginally above this threshold ($d = 0.502$).

Physically, this makes sense: tankers have a higher beam-to-length ratio than
cargo vessels of comparable size. When viewed broadside at typical coastal ranges
($r = 5$--$20$~km), a tanker subtends a larger azimuth angle than a cargo ship,
producing a measurably larger \texttt{az\_extent\_m}. However, because the
aspect angle varies with vessel orientation, this signal is noisy —
a cargo vessel viewed broadside may appear wider than a tanker viewed bow-on.

\subsection{Classical ML vs.\ Quantum ML}

\textit{This section will be completed with final results from all models.}

Key dimensions of comparison:
\begin{itemize}
    \item \textbf{Accuracy and F1}: Do QML models match or exceed GBT?
    \item \textbf{Confidence calibration}: Are QML confidence scores better
          separated between correct/incorrect predictions?
    \item \textbf{Training efficiency}: How do epoch counts and convergence
          speed compare?
    \item \textbf{Data efficiency}: Do QML models achieve competitive
          performance with fewer training samples?
\end{itemize}

\subsection{Limitations}

\begin{enumerate}
    \item \textbf{Quantum simulation}: All QML experiments are performed on
          classical quantum simulators. Simulation overhead scales exponentially
          with qubit count; results on real quantum hardware may differ.
    \item \textbf{Feature ceiling}: With all Cohen's $d < 0.8$ for the 70/80
          boundary, there is a fundamental limit to classification accuracy
          achievable from single-snapshot radar features alone.
    \item \textbf{Temporal features}: Rolling statistics (SOG, COG variability)
          were not used as primary features in the initial experiments due to
          high null rates in short tracks. Their inclusion may further improve
          Cargo/Tanker separation.
    \item \textbf{Label quality}: AIS type codes are self-reported and may
          contain errors. Label noise acts as a hard ceiling on achievable accuracy.
\end{enumerate}

%% ============================================================
\section{Conclusion}
\label{sec:conclusion}

This paper presents the first systematic study of multi-class vessel type
classification from terrestrial coastal radar using both classical and
quantum machine learning. We have shown that:
\begin{enumerate}
    \item The physical derivation of radar size features matters significantly:
          computing the cross-range extent in physical metres (\texttt{az\_extent\_m})
          is essential for separating Cargo from Tanker, whereas the raw azimuth
          gate width in radians produces near-zero ellipse areas and degenerate
          Euclidean sizes.
    \item \textit{(Further conclusions pending final model comparison.)}
\end{enumerate}

\section*{Acknowledgements}
\textit{[To be completed]}

%% ============================================================
\bibliographystyle{plainnat}
\begin{thebibliography}{99}

\bibitem{breiman2001random}
Breiman, L. (2001).
\newblock Random forests.
\newblock \textit{Machine Learning}, 45(1), 5--32.

\bibitem{chawla2002smote}
Chawla, N.~V., Bowyer, K.~W., Hall, L.~O., \& Kegelmeyer, W.~P. (2002).
\newblock SMOTE: Synthetic minority over-sampling technique.
\newblock \textit{Journal of Artificial Intelligence Research}, 16, 321--357.

\bibitem{cortes1995support}
Cortes, C., \& Vapnik, V. (1995).
\newblock Support-vector networks.
\newblock \textit{Machine Learning}, 20(3), 273--297.

\bibitem{friedman2001greedy}
Friedman, J.~H. (2001).
\newblock Greedy function approximation: A gradient boosting machine.
\newblock \textit{Annals of Statistics}, 29(5), 1189--1232.

\bibitem{havlicek2019supervised}
Havl{\'i}{\v{c}}ek, V., C{\'o}rcoles, A.~D., Temme, K., Harrow, A.~W.,
Kandala, A., Chow, J.~M., \& Gambetta, J.~M. (2019).
\newblock Supervised learning with quantum-enhanced feature spaces.
\newblock \textit{Nature}, 567(7747), 209--212.

\bibitem{mari2020transfer}
Mari, A., Bromley, T.~R., Izaac, J., Schuld, M., \& Killoran, N. (2020).
\newblock Transfer learning in hybrid classical-quantum neural networks.
\newblock \textit{Quantum}, 4, 340.

\bibitem{pennylane2019}
Bergholm, V., Izaac, J., Schuld, M., \textit{et al.} (2018).
\newblock PennyLane: Automatic differentiation of hybrid quantum-classical computations.
\newblock \textit{arXiv preprint arXiv:1811.04975}.

\bibitem{schuld2020circuit}
Schuld, M., Bocharov, A., Svore, K.~M., \& Wiebe, N. (2020).
\newblock Circuit-centric quantum classifiers.
\newblock \textit{Physical Review A}, 101(3), 032308.

\end{thebibliography}

\end{document}
